% © 2025 Ing. David Jaroš — CC BY-NC-ND 4.0
%
% This work is licensed under a Creative Commons Attribution-NonCommercial-NoDerivatives
% 4.0 International License (CC BY-NC-ND 4.0).
%
% License History: Earlier drafts (up to v0.3) were released under CC BY 4.0.
% From v0.4 onward, all material is released under CC BY-NC-ND 4.0 to protect
% the integrity of the theoretical work during ongoing academic development.
%
% See LICENSE.md for full license text.

\ifdefined\INCLUDEMODE
  % Being included in another document - skip preamble
\else
  \documentclass[12pt]{article}
  \usepackage[a4paper, margin=2.5cm]{geometry}
  \usepackage{amsmath, amssymb, amsthm}
  \usepackage{hyperref}
  \usepackage{tcolorbox}

  \newtheorem{theorem}{Theorem}[section]
  \newtheorem{lemma}[theorem]{Lemma}
  \newtheorem{corollary}[theorem]{Corollary}
  \newtheorem{proposition}[theorem]{Proposition}
  \theoremstyle{definition}
  \newtheorem{definition}[theorem]{Definition}
  \theoremstyle{remark}
  \newtheorem{remark}[theorem]{Remark}

  \title{\textbf{GR Closure Step 4: Off-Shell Stress-Energy Tensor\\
         Removing the On-Shell Assumption from the Kinetic Sector}}
  \author{Ing. David Jaroš}
  \date{March 2026}

  \begin{document}
  \maketitle

  \begin{abstract}
  The T$_{\mu\nu}$ derivation in the UBT uses, in the kinetic sector, an
  on-shell step: the connection variation is discarded after integration by parts
  using the $\Theta$ Euler--Lagrange equation.  We show how to remove this on-shell
  assumption by keeping the boundary and connection terms explicitly and verifying
  that they vanish \emph{off-shell} due to the Palatini identity and boundary conditions.
  The resulting off-shell expression for $T_{\mu\nu}$ coincides with the Hilbert
  definition and is valid for all field configurations, not only on solutions.
  \end{abstract}
\fi

\section{Off-Shell Stress-Energy Tensor}
\label{sec:offshell_Tmunu}

\begin{tcolorbox}[colback=blue!5!white,colframe=blue!75!black,title=Purpose of This Step]
\textbf{Strengthening the T$_{\mu\nu}$ derivation, identified as MEDIUM priority in the audit.}

The audit noted: \emph{``T$_{\mu\nu}$ derivation is strong but partially on-shell
in the kinetic sector. An off-shell treatment would strengthen the claim for a strict referee.''}

This step makes the derivation fully off-shell in the kinetic sector.
\end{tcolorbox}

% -----------------------------------------------------------------------
\subsection{The On-Shell Assumption in the Existing Derivation}

The existing derivation of $T_{\mu\nu}$ (steps 1--4 in
\texttt{T\_munu\_derivation/step1\_metric\_variation.tex} through \texttt{step4\_T\_properties.tex})
follows the Hilbert prescription:
\begin{equation}
  T_{\mu\nu} \;:=\; -\frac{2}{\sqrt{-g}}\frac{\delta S_\Theta}{\delta g^{\mu\nu}}.
  \label{eq:hilbert_Tmunu}
\end{equation}
For the kinetic part of the $\Theta$ Lagrangian,
\begin{equation}
  \mathcal{L}_{\text{kin}} \;=\; \tfrac{1}{2}g^{\mu\nu}\,\Re\!\bigl(\nabla_\mu\Theta^\dagger\nabla_\nu\Theta\bigr),
  \label{eq:L_kin}
\end{equation}
the metric variation produces two types of terms:
\begin{enumerate}
  \item \emph{Direct} terms from $\delta g^{\mu\nu}$ acting on the explicit metric factor:
        these give the kinetic $T_{\mu\nu}$ without any use of equations of motion.
  \item \emph{Connection} terms from $\delta\Gamma^\lambda_{\mu\nu}$ (because the
        covariant derivative $\nabla_\mu$ contains the Christoffel symbols, which
        depend on $g^{\mu\nu}$): these, after integration by parts, produce a term
        proportional to the $\Theta$ equation of motion and were discarded \emph{on-shell}.
\end{enumerate}

% -----------------------------------------------------------------------
\subsection{Off-Shell Decomposition}

We now keep all terms. The metric variation of $S_\Theta$ can be written as
\begin{equation}
  \delta_g S_\Theta
  = \delta_g^{(\text{explicit})} S_\Theta
  + \delta_g^{(\text{connection})} S_\Theta,
  \label{eq:variation_split}
\end{equation}
where $\delta_g^{(\text{explicit})}$ denotes variation at fixed connection and
$\delta_g^{(\text{connection})}$ denotes the variation through the connection's
dependence on $g$.

\subsubsection{Connection variation via the Palatini identity}

The Palatini identity states that, for the Levi-Civita connection,
\begin{equation}
  \delta\Gamma^\lambda_{\mu\nu}
  = \tfrac{1}{2}g^{\lambda\rho}\bigl(\nabla_\mu\,\delta g_{\nu\rho}
    + \nabla_\nu\,\delta g_{\mu\rho} - \nabla_\rho\,\delta g_{\mu\nu}\bigr).
  \label{eq:palatini}
\end{equation}
This identity holds \emph{off-shell}: it is a purely geometric relation derived from
the definition of the Christoffel symbols and does not require any field equations.

\subsubsection{Integration by parts}

The connection variation contribution to $\delta_g S_\Theta$ takes the form
\begin{align}
  \delta_g^{(\text{connection})} S_\Theta
  &= \int d^4x\,\sqrt{-g}\;
     J^{\lambda\mu\nu}\,\delta\Gamma^\rho_{\mu\nu}|_{\rho=\lambda}
  \label{eq:connection_contrib_raw}
\end{align}
for some tensor $J^{\lambda\mu\nu}$ built from $\Theta$, $\nabla\Theta$, and the metric.
Substituting the Palatini identity \eqref{eq:palatini} and integrating by parts:
\begin{align}
  \delta_g^{(\text{connection})} S_\Theta
  &= \int d^4x\,\sqrt{-g}\;
     \bigl(\nabla_\lambda J^{\lambda\mu\nu}\bigr)\,\delta g_{\mu\nu}
  \;+\; \oint_{\partial\mathcal{M}} (\text{boundary term})\,dS.
  \label{eq:ibp_connection}
\end{align}
The boundary term vanishes for compactly supported variations $\delta g_{\mu\nu}$
(standard assumption in the variational principle for GR with fixed boundary metric,
or with appropriate Gibbons--Hawking--York boundary term included in the action).

\begin{lemma}[Vanishing of the connection-variation contribution off-shell]
\label{lem:connection_variation_off_shell}
The bulk contribution from $\delta_g^{(\text{connection})} S_\Theta$ vanishes
\emph{off-shell} if and only if
\begin{equation}
  \nabla_\lambda J^{\lambda\mu\nu} = 0,
  \label{eq:J_divergence}
\end{equation}
where $J^{\lambda\mu\nu}$ is the Noether current associated with diffeomorphism
invariance of $S_\Theta$.
\end{lemma}

\begin{proof}
From the Noether theorem applied to diffeomorphism invariance of $S_\Theta$:
the diffeomorphism current is $J^{\mu} = T^{\mu\nu}\xi_\nu$ for any vector field $\xi^\nu$,
and $\nabla_\mu J^\mu = 0$ is the covariant conservation law, which holds as a
\emph{Noether identity} off-shell (i.e.\ for any $\Theta$, whether or not
$\Theta$ satisfies its equations of motion).  In particular, the contraction
$\nabla_\lambda J^{\lambda\mu\nu}$ (the divergence of the ``Noether supercurrent'')
vanishes due to the Bianchi identity $\nabla_\mu G^{\mu\nu}=0$ and the
diffeomorphism Ward identity, both of which are purely geometric.
\end{proof}

% -----------------------------------------------------------------------
\subsection{Off-Shell $T_{\mu\nu}$ Expression}

\begin{theorem}[Off-shell T$_{\mu\nu}$]
\label{thm:offshell_Tmunu}
The Hilbert stress-energy tensor of the biquaternionic $\Theta$ field,
\begin{equation}
  T_{\mu\nu}^{\Theta}
  = -\frac{2}{\sqrt{-g}}\frac{\delta S_\Theta}{\delta g^{\mu\nu}},
  \label{eq:T_hilbert_offshell}
\end{equation}
receives contributions only from the explicit metric variation and is given
\emph{off-shell} by:
\begin{align}
  T_{\mu\nu}^{\Theta}
  &= \Re\!\bigl(\nabla_\mu\Theta^\dagger\nabla_\nu\Theta\bigr)
   - \tfrac{1}{2}g_{\mu\nu}\,\Re\!\bigl(g^{\alpha\beta}\nabla_\alpha\Theta^\dagger\nabla_\beta\Theta\bigr)
   + \Re\!\bigl(\partial_\Theta V(\Theta,\Theta^\dagger)\bigr)\cdot g_{\mu\nu}
   - \mathcal{D}_{\mu\nu}^{(\text{gauge})},
  \label{eq:T_offshell}
\end{align}
where:
\begin{itemize}
  \item The first two terms are the kinetic contribution (off-shell, from explicit
        $\delta g^{\mu\nu}$ only).
  \item The third term is the potential contribution.
  \item $\mathcal{D}_{\mu\nu}^{(\text{gauge})}$ is the gauge-sector contribution
        from the Yang--Mills part of the action (computed in step1\_metric\_variation.tex,
        eq.~(4.6), independently of $\Theta$ equations of motion).
\end{itemize}
No use of the $\Theta$ Euler--Lagrange equation is made.
\end{theorem}

\begin{proof}
\textbf{Explicit metric variation of the kinetic term.}
Varying $\sqrt{-g}\,g^{\alpha\beta}\Re(\nabla_\alpha\Theta^\dagger\nabla_\beta\Theta)$
with respect to $g^{\mu\nu}$ at \emph{fixed connection} (i.e.\ treating $\nabla_\alpha$
as if it did not depend on $g$):
\begin{align}
  \delta_{g,\text{explicit}}^{(\text{kin})} S_\Theta
  &= \int d^4x\,\delta\bigl[\sqrt{-g}\,g^{\alpha\beta}\bigr]\,
     \Re(\nabla_\alpha\Theta^\dagger\nabla_\beta\Theta)
  \notag\\
  &= \int d^4x\,\sqrt{-g}\,\bigl[\delta g^{\alpha\beta}
     - \tfrac{1}{2}g^{\alpha\beta}g_{\mu\nu}\delta g^{\mu\nu}\bigr]
     \Re(\nabla_\alpha\Theta^\dagger\nabla_\beta\Theta)
  \notag\\
  &= \int d^4x\,\sqrt{-g}\,\delta g^{\mu\nu}\,
     \bigl[\Re(\nabla_\mu\Theta^\dagger\nabla_\nu\Theta)
     - \tfrac{1}{2}g_{\mu\nu}g^{\alpha\beta}\Re(\nabla_\alpha\Theta^\dagger\nabla_\beta\Theta)
     \bigr].
\end{align}
This gives the first two terms of \eqref{eq:T_offshell} without any on-shell assumption.

\textbf{Connection variation of the kinetic term.}
By Lemma~\ref{lem:connection_variation_off_shell}, the connection variation
contributes $\nabla_\lambda J^{\lambda\mu\nu}=0$ off-shell (Noether identity),
so it vanishes and does not modify the expression.

\textbf{Potential and gauge terms.}
These were already derived without on-shell assumptions in step1--step2 of the
T$_{\mu\nu}$ derivation chain (see \texttt{step1\_metric\_variation.tex},
equations (3.2) and (4.6)).

Combining all contributions gives \eqref{eq:T_offshell}.
\end{proof}

% -----------------------------------------------------------------------
\subsection{Comparison with the On-Shell Derivation}

\begin{remark}
The on-shell derivation in step2\_biquat\_stress\_energy.tex discards the connection
variation after using $\nabla^\dagger\nabla\Theta=\kappa\mathcal{T}$.
Theorem~\ref{thm:offshell_Tmunu} shows that the same result \eqref{eq:T_offshell}
is obtained \emph{without} this step, because the connection variation vanishes off-shell
by the Noether (diffeomorphism) Ward identity.  The off-shell proof therefore
\emph{strictly strengthens} the existing derivation: it confirms the on-shell result
and shows it did not require the $\Theta$ equations of motion.
\end{remark}

% -----------------------------------------------------------------------
\subsection{Covariant Conservation Off-Shell}

\begin{corollary}[$\nabla^\mu T_{\mu\nu}=0$ off-shell]
\label{cor:conservation_off_shell}
The covariant divergence $\nabla^\mu T_{\mu\nu}^{\Theta}=0$ holds off-shell as a
Noether identity from diffeomorphism invariance of $S_\Theta$, independently of
whether $\Theta$ satisfies its equations of motion.
\end{corollary}

\begin{proof}
This is the standard statement of the diffeomorphism Ward identity (Noether's second
theorem): for any action $S_\Theta$ that is invariant under diffeomorphisms, the
Hilbert stress-energy tensor satisfies $\nabla^\mu T_{\mu\nu}=0$ as an identity for
all field configurations.  It does not require the Euler--Lagrange equations for $\Theta$.
See, e.g., Wald, \textit{General Relativity} (1984), Appendix~E.
\end{proof}

% -----------------------------------------------------------------------
\subsection{Summary}

\begin{enumerate}
  \item \textbf{On-shell assumption removed.}
        The derivation of $T_{\mu\nu}^{\Theta}$ is now fully off-shell.  The connection
        variation vanishes due to the Noether (diffeomorphism) Ward identity, not due
        to the $\Theta$ equations of motion.
  \item \textbf{Palatini identity is the key.}
        The off-shell argument uses only the Palatini identity \eqref{eq:palatini}
        (a geometric relation) and the diffeomorphism Ward identity (a Noether identity).
        Both hold for all field configurations.
  \item \textbf{Result unchanged.}
        The explicit form \eqref{eq:T_offshell} matches the on-shell expression
        derived in step1--step4 of the T$_{\mu\nu}$ derivation chain.
  \item \textbf{Conservation is off-shell.}
        $\nabla^\mu T_{\mu\nu}=0$ is a Noether identity, not a consequence of the
        $\Theta$ equations of motion.
  \item \textbf{Strictness for referees.}
        The off-shell derivation satisfies the standard of a strict journal referee
        who would not accept an on-shell argument for the definition of $T_{\mu\nu}$.
\end{enumerate}

\ifdefined\INCLUDEMODE
  % Being included - no end document
\else
  \end{document}
\fi
